\documentclass[10pt,a4paper]{amsart}
\usepackage{amsmath,amssymb,amsthm,graphicx}
% Retained preamble for the Stacks Project chapter snapshot in this directory.
% The chapter's own preamble.tex is a copy of the shared Stacks Project preamble; the block
% below is sliced MECHANICALLY, line for line, from it. Nothing was added, removed or
% redefined; the wrapper class, the xr-hyper/hyperref wiring and the externaldocument list
% are omitted because the assembler supplies the document class.
%
% lines taken from preamble.tex: 3, 5, 9, 11, 12, 156, 157, 158, ... (51 lines total)
%
% package coverage restored from the chapter's own preamble (2026-10-13): xy, xy-2cell

\usepackage[all]{xy}
\xyoption{2cell}
\UseAllTwocells

\newenvironment{reference}{\comment}{\endcomment}
\newenvironment{slogan}{\comment}{\endcomment}
\newenvironment{history}{\comment}{\endcomment}
\theoremstyle{plain}
\newtheorem{theorem}[subsection]{Theorem}
\newtheorem{proposition}[subsection]{Proposition}
\newtheorem{lemma}[subsection]{Lemma}
\theoremstyle{definition}
\newtheorem{definition}[subsection]{Definition}
\newtheorem{example}[subsection]{Example}
\newtheorem{exercise}[subsection]{Exercise}
\newtheorem{situation}[subsection]{Situation}
\theoremstyle{remark}
\newtheorem{remark}[subsection]{Remark}
\newtheorem{remarks}[subsection]{Remarks}
\numberwithin{equation}{subsection}
\def\lim{\mathop{\mathrm{lim}}\nolimits}
\def\colim{\mathop{\mathrm{colim}}\nolimits}
\def\Spec{\mathop{\mathrm{Spec}}}
\def\Hom{\mathop{\mathrm{Hom}}\nolimits}
\def\Ext{\mathop{\mathrm{Ext}}\nolimits}
\def\SheafHom{\mathop{\mathcal{H}\!\mathit{om}}\nolimits}
\def\SheafExt{\mathop{\mathcal{E}\!\mathit{xt}}\nolimits}
\def\Sch{\mathit{Sch}}
\def\Mor{\mathop{\mathrm{Mor}}\nolimits}
\def\Ob{\mathop{\mathrm{Ob}}\nolimits}
\def\Sh{\mathop{\mathit{Sh}}\nolimits}
\def\NL{\mathop{N\!L}\nolimits}
\def\CH{\mathop{\mathrm{CH}}\nolimits}
\def\proetale{{pro\text{-}\acute{e}tale}}
\def\etale{{\acute{e}tale}}
\def\QCoh{\mathit{QCoh}}
\def\Ker{\mathop{\mathrm{Ker}}}
\def\Im{\mathop{\mathrm{Im}}}
\def\Coker{\mathop{\mathrm{Coker}}}
\def\Coim{\mathop{\mathrm{Coim}}}
\DeclareMathSymbol{\boxtimes}{\mathbin}{AMSa}{"02}
\def\QCohstack{\mathcal{QC}\!\mathit{oh}}
\def\Cohstack{\mathcal{C}\!\mathit{oh}}
\def\Spacesstack{\mathcal{S}\!\mathit{paces}}
\def\Quotfunctor{\mathrm{Quot}}
\def\Hilbfunctor{\mathrm{Hilb}}
\def\Curvesstack{\mathcal{C}\!\mathit{urves}}
\def\Polarizedstack{\mathcal{P}\!\mathit{olarized}}
\def\Complexesstack{\mathcal{C}\!\mathit{omplexes}}
\def\Pic{\mathop{\mathrm{Pic}}\nolimits}
\def\Picardstack{\mathcal{P}\!\mathit{ic}}
\def\Picardfunctor{\mathrm{Pic}}
\def\Deformationcategory{\mathcal{D}\!\mathit{ef}}
\usepackage{cleveref}
\crefname{originallemma}{original lemma}{original lemmas}
\Crefname{originallemma}{Original lemma}{Original lemmas}
\makeatletter
\newlabel{duality-section-duality-finite}{{0AWZ}{source}{}{sourcepointer-1}{}}
\newlabel{duality-section-duality-finite@cref}{{[originallemma][0][]0AWZ}{[1][0][]source}{}{}{}}
\makeatother
\crefname{originallemma}{original lemma}{original lemmas}
\Crefname{originallemma}{Original lemma}{Original lemmas}
\makeatletter
\newlabel{duality-section-upper-shriek}{{0A9Y}{source}{}{sourcepointer-2}{}}
\newlabel{duality-section-upper-shriek@cref}{{[originallemma][0][]0A9Y}{[1][0][]source}{}{}{}}
\makeatother
\crefname{originallemma}{original lemma}{original lemmas}
\Crefname{originallemma}{Original lemma}{Original lemmas}
\makeatletter
\newlabel{lemma-compare-dualizing}{{0BUL}{source}{}{sourcepointer-3}{}}
\newlabel{lemma-compare-dualizing@cref}{{[originallemma][0][]0BUL}{[1][0][]source}{}{}{}}
\makeatother
\crefname{originallemma}{original lemma}{original lemmas}
\Crefname{originallemma}{Original lemma}{Original lemmas}
\makeatletter
\newlabel{lemma-finite-flat-trace}{{0BT8}{source}{}{sourcepointer-4}{}}
\newlabel{lemma-finite-flat-trace@cref}{{[originallemma][0][]0BT8}{[1][0][]source}{}{}{}}
\makeatother
\crefname{originallemma}{original lemma}{original lemmas}
\Crefname{originallemma}{Original lemma}{Original lemmas}
\makeatletter
\newlabel{remark-relative-dualizing-for-flat-quasi-finite}{{0BVJ}{source}{}{sourcepointer-5}{}}
\newlabel{remark-relative-dualizing-for-flat-quasi-finite@cref}{{[originallemma][0][]0BVJ}{[1][0][]source}{}{}{}}
\makeatother
\crefname{originallemma}{original lemma}{original lemmas}
\Crefname{originallemma}{Original lemma}{Original lemmas}
\makeatletter
\newlabel{section-discriminant}{{0BVH}{source}{}{sourcepointer-6}{}}
\newlabel{section-discriminant@cref}{{[originallemma][0][]0BVH}{[1][0][]source}{}{}{}}
\makeatother
\title{Stacks, Discriminants: lemma-compare-trace}
\author{The Stacks Project Authors}
\date{stacks-project:discriminant.tex:local-snapshot-20261012; GNU FDL 1.2 or later}
\begin{document}
\maketitle
\noindent\textit{Editorial note.} Statement span is lines 3433-3447 of sources/stacks-discriminant/pinned-original-source.tex; proof span is lines 3449-3492, the complete author proof (42 body lines).\par\medskip
\noindent\textit{Editorial note.} Navigation stubs: 6 cited-but-not-included labels are supplied as non-load-bearing pointers with their pinned tags where the snapshot tag index has them.\par\medskip
\noindent\textit{Editorial note.} extra\_packages is empty on purpose: the retained preamble already loads xy with options and a second declaration clashes in the assembler wrapper.\par\medskip
\noindent\textit{Editorial note.} Text is retained verbatim, including the author's own typos and placeholders.\par\medskip
\noindent\textit{Editorial note.} Queue index 9 of queue\_discriminant.json (pb=1777, stmt 3433-3447, proof 3449-3492).\par\medskip
\noindent\textit{Allowed standard prerequisites (editorial).}\begin{itemize}
\item Explicit prerequisite boundary. The cited inputs are prerequisites; the argument in the retained span is not replaced by any cited result.
\end{itemize}\medskip
\ifdefined\hypertarget\hypertarget{sourcepointer-1}{}\fi
\noindent\textit{Source reference.} Original lemma 0AWZ: sources/stacks-discriminant/pinned-original-source.tex, source lines 3492--3492. Cited by the retained statement or proof but not part of the collected spans; supplied only as a navigation pointer. The argument does not depend on its content.\par\medskip
\ifdefined\hypertarget\hypertarget{sourcepointer-2}{}\fi
\noindent\textit{Source reference.} Original lemma 0A9Y: sources/stacks-discriminant/pinned-original-source.tex, source lines 3492--3492. Cited by the retained statement or proof but not part of the collected spans; supplied only as a navigation pointer. The argument does not depend on its content.\par\medskip
\ifdefined\hypertarget\hypertarget{sourcepointer-3}{}\fi
\noindent\textit{Source reference.} Original lemma 0BUL: sources/stacks-discriminant/pinned-original-source.tex, source lines 3360--3375. Cited by the retained statement or proof but not part of the collected spans; supplied only as a navigation pointer. The argument does not depend on its content.\par\medskip
\ifdefined\hypertarget\hypertarget{sourcepointer-4}{}\fi
\noindent\textit{Source reference.} Original lemma 0BT8: sources/stacks-discriminant/pinned-original-source.tex, source lines 874--877. Cited by the retained statement or proof but not part of the collected spans; supplied only as a navigation pointer. The argument does not depend on its content.\par\medskip
\ifdefined\hypertarget\hypertarget{sourcepointer-5}{}\fi
\noindent\textit{Source reference.} Original lemma 0BVJ: sources/stacks-discriminant/pinned-original-source.tex, source lines 1107--1124. Cited by the retained statement or proof but not part of the collected spans; supplied only as a navigation pointer. The argument does not depend on its content.\par\medskip
\ifdefined\hypertarget\hypertarget{sourcepointer-6}{}\fi
\noindent\textit{Source reference.} Original lemma 0BVH: sources/stacks-discriminant/pinned-original-source.tex, source lines 694--759. Cited by the retained statement or proof but not part of the collected spans; supplied only as a navigation pointer. The argument does not depend on its content.\par\medskip
\begin{lemma}
\label{lemma-compare-trace}
Let $f : Y \to X$ be a finite flat morphism of Noetherian schemes.
The map
$$
\text{Trace}_f : f_*\mathcal{O}_Y \longrightarrow \mathcal{O}_X
$$
of Section \ref{section-discriminant}
corresponds to a map $\mathcal{O}_Y \to f^!\mathcal{O}_X$ (see proof).
Denote $\tau_{Y/X} \in H^0(Y, f^!\mathcal{O}_X)$ the image of $1$.
Via the isomorphism $H^0(f^!\mathcal{O}_X) = \omega_{X/Y}$ of
Lemma \ref{lemma-compare-dualizing}
this agrees with the construction in
Remark \ref{remark-relative-dualizing-for-flat-quasi-finite}.
\end{lemma}

\begin{proof}
The functor $f^!$ is defined in 
Duality for Schemes, Section \ref{duality-section-upper-shriek}.
Since $f$ is finite (and hence proper), we see that $f^!$ is given by
the right adjoint to pushforward for $f$. In
Duality for Schemes, Section \ref{duality-section-duality-finite}
we have made this adjoint explicit. In particular,
the object $f^!\mathcal{O}_X$ consists of a single
cohomology sheaf placed in degree $0$ and for this sheaf we have
$$
f_*f^!\mathcal{O}_X =
\SheafHom_{\mathcal{O}_X}(f_*\mathcal{O}_Y, \mathcal{O}_X)
$$
To see this we use also that $f_*\mathcal{O}_Y$ is finite locally free
as $f$ is a finite flat morphism of Noetherian schemes
and hence all higher Ext sheaves are zero. Some details omitted.
Thus finally
$$
\text{Trace}_f \in
\Hom_{\mathcal{O}_X}(f_*\mathcal{O}_Y, \mathcal{O}_X) =
\Gamma(X, f_*f^!\mathcal{O}_X) =
\Gamma(Y, f^!\mathcal{O}_X)
$$
On the other hand, we have $f^!\mathcal{O}_X = \omega_{Y/X}$
by the identification of Lemma \ref{lemma-compare-dualizing}.
Thus we now have two elements, namely $\text{Trace}_f$
and $\tau_{Y/X}$ from
Remark \ref{remark-relative-dualizing-for-flat-quasi-finite} in
$$
\Gamma(Y, f^!\mathcal{O}_X) = \Gamma(Y, \omega_{Y/X})
$$
and the lemma says these elements are the same.

\medskip\noindent
Let $U = \Spec(A) \subset X$ be an affine open with inverse image
$V = \Spec(B) \subset Y$. Since $f$ is finite, we see that
$A \to B$ is finite and hence the $\omega_{Y/X}(V) = \Hom_A(B,A)$
by construction and this isomorphism agrees with the identification
of $f_*f^!\mathcal{O}_Y$ with
$\SheafHom_{\mathcal{O}_X}(f_*\mathcal{O}_Y, \mathcal{O}_X)$ discussed
above. Hence the agreement of $\text{Trace}_f$ and $\tau_{Y/X}$
follows from the fact that $\tau_{B/A} = \text{Trace}_{B/A}$
by Lemma \ref{lemma-finite-flat-trace}.
\end{proof}

\end{document}
